\documentclass[../../../include/open-logic-section]{subfiles}

\begin{document}

\olfileid[hi]{sth}{choice}{intro}

\olsection{परिचय}

\crefrange{sth:cardinals::chap}{sth:card-arithmetic::chap} में हमने
कार्डिनलों को क्रमसूचक मानकर कार्डिनलों का सिद्धांत विकसित किया। यह
दृष्टिकोण सुव्यवस्था अभिगृहीत पर निर्भर है। पता चलता है कि सुव्यवस्था एक
अन्य सिद्धांत---चयन अभिगृहीत---के तुल्य है और उसकी स्वीकार्यता पर गंभीर
दार्शनिक चर्चा हुई है। इस अध्याय में हमारे प्रश्न हैं: इस अभिगृहीत का उपयोग
कैसे होता है और क्या उसका औचित्य दिया जा सकता है?

\end{document}
